\documentclass[11pt]{article}

\usepackage[margin=1in]{geometry}
\usepackage{booktabs}
\usepackage{amsmath}
\usepackage{microtype}
\usepackage[T1]{fontenc}
\usepackage{ebgaramond}
\usepackage[hidelinks]{hyperref}
\usepackage{xcolor}
\definecolor{oxblood}{HTML}{7A2424}
\hypersetup{colorlinks=true,urlcolor=oxblood,linkcolor=oxblood}

\setlength{\parskip}{0.4em}
\setlength{\parindent}{0pt}

\title{\textbf{Prompt style, not translation register,\\drives autoformalization of Euclid into Lean 4}}
\author{Reza Ramji}
\date{}

\begin{document}
\maketitle

\begin{abstract}
\noindent A proof in a textbook is written for people; Lean~4 is a language a
computer can verify line by line. Autoformalization is the task of translating
ordinary mathematical prose into Lean automatically. Using Euclid's
\emph{Elements} Book~I and GPT-4o, we ask whether the English wording of a
proposition changes how well it formalizes. A counterbalanced $3\times3$ design
(three translation registers $\times$ three prompt styles, 2{,}653 runs) shows
that prompt style is the lever and translation register is not: prompt style
spans a 24-point range in compile rate ($p<0.001$), while translation has no
significant effect ($p=0.603$). An earlier ``one translation wins'' result was
a prompt--translation confound. A separate probe shows a small local model had
memorized the benchmark, forcing the switch to GPT-4o.
\end{abstract}

\section{The question}

Lean~4 is a proof language a machine can check: if a Lean proof compiles and
type-checks, it is correct, with no appeal to trust. Translating informal
mathematics into Lean by hand is slow expert work, so the hope is that a
language model can do it---``autoformalization.'' We test a narrow version of
that hope on Euclid's \emph{Elements}, built on the
\href{https://github.com/loganrjmurphy/LeanEuclid}{LeanEuclid} benchmark
(Murphy et al., 2024).

Book~I---the same 48 propositions---exists in sharply different English
registers. We use three: \textbf{Heath} (1908, archaic Victorian),
\textbf{Fitzpatrick} (2007, modern academic with variables named explicitly),
and a plain \textbf{Modern} paraphrase. All three say the same things. Does the
register change how well GPT-4o formalizes the proposition?

\section{The confound, and the fix}

An early run used a single prompt template whose five worked examples were all
in Fitzpatrick's style. It appeared that Fitzpatrick text formalized best,
beating Heath by 8.4~percentage points ($p=0.008$). But the Fitzpatrick
translation matched the \emph{examples in the prompt}, so register and
prompt-style matching were confounded.

To separate them we crossed the two factors: three translation registers
against three prompt styles (Fitzpatrick-, Heath-, and Modern-style examples),
a fully crossed $3\times3$ design over 143 propositions (Book~I plus 100 UniGeo
theorems), for 2{,}653 total runs.

\section{Result: prompt style dominates}

How you prompt matters; which translation you feed in does not. Prompt style
swings the compile rate across a \textbf{24-point range}---72.6\% with the best
style versus 48.4\% with the worst ($p<0.001$). Translation register has no
significant effect once prompt style is held fixed ($p=0.603$).

\begin{table}[h]
\centering
\caption*{\textit{Compile rate by translation $\times$ prompt style (Book~I, GPT-4o)}}
\begin{tabular}{lccc}
\toprule
 & Fitzpatrick prompt & Heath prompt & Modern prompt \\
\midrule
Fitzpatrick & \textbf{76.3\%} & 60.0\% & 45.9\% \\
Heath       & 67.9\%          & 65.5\% & 52.6\% \\
Modern      & 73.5\%          & 63.2\% & 46.8\% \\
\bottomrule
\end{tabular}
\end{table}

Read down each column (prompt fixed) and the translations barely separate; read
across the columns (prompt changing) and the rates collapse. The original
``Fitzpatrick wins'' gap appears only in the Fitzpatrick-prompt column; under
Heath-style prompts it reverses (Heath 65.5\% $>$ Fitzpatrick 60.0\%), and
under Modern-style prompts it reverses again (52.6\% $>$ 45.9\%). On the 100
independent UniGeo theorems, Heath beat Fitzpatrick by 6.2pp---the opposite of
the first claim.

\section{Statistical model}

A mixed-effects linear regression with proposition as a random intercept
($N=2{,}653$) confirms the pattern: translation coefficients are tiny and
non-significant, prompt coefficients large and significant.

\begin{table}[h]
\centering
\caption*{\textit{Mixed-effects model coefficients}}
\begin{tabular}{lccc}
\toprule
Predictor & Coefficient & SE & $p$ \\
\midrule
Translation (Fitzpatrick vs.\ Heath) & $+0.010$ & 0.019 & 0.603 \\
Translation (Modern vs.\ Heath)      & $+0.009$ & 0.021 & 0.665 \\
Prompt (Heath vs.\ Fitzpatrick)      & $\mathbf{-0.097}$ & 0.022 & $\mathbf{<0.001}$ \\
Prompt (Modern vs.\ Fitzpatrick)     & $\mathbf{-0.242}$ & 0.024 & $\mathbf{<0.001}$ \\
\bottomrule
\end{tabular}
\end{table}

Adding prompt style improves model fit by 99.5 AIC points; adding a
prompt$\times$translation interaction does not help (AIC rises by 4.3). A paired
$t$-test on per-proposition rates ($N=66$ items, marginalized over prompt)
leaves the two translations indistinguishable: Heath 70.0\% vs.\ Fitzpatrick
62.3\%, $t(65)=-1.66$, $p=0.102$, Cohen's $d=-0.20$.

\section{The contamination caveat}

We first tried a small local model, Qwen2.5-Coder-7B (Q4\_K\_M), on a Jetson
Orin Nano Super (8\,GB, ${\sim}9.5$ tokens/s). It had memorized the benchmark.
Its \emph{best} compile rate came from no math at all---just ``Proposition N.''

\begin{table}[h]
\centering
\caption*{\textit{Contamination probe --- Qwen2.5-Coder-7B}}
\begin{tabular}{lcc}
\toprule
Input given & Compile rate & LeanEuclid vocab \\
\midrule
Full Heath text          & 60.5\% & 100\% \\
Empty (``Proposition N'') & \textbf{62.8\%} & 100\% \\
Garbled vocabulary       & 60.5\% & 100\% \\
Shuffled clauses         & 60.5\% & 100\% \\
\bottomrule
\end{tabular}
\end{table}

The highest score from empty input means the model recites memorized structure
rather than reading the math. Semantic equivalence does drop to 0\% on empty
input (versus 16--27\% with garbled text), so it reads input to refine meaning
but does not need it to produce something that compiles. GPT-4o, by contrast,
is clean: 10\% from empty input versus 67.9\% from full text. We dropped Qwen
and ran everything counterbalanced on GPT-4o.

\section{Takeaway}

If you want to use autoformalization on your own mathematics, the lever is not
cleaning up your prose---it is styling your few-shot examples to match how you
write, especially whether they introduce variables explicitly. Two
methodological notes follow: a single prompt template can manufacture a spurious
``wording'' effect, so counterbalanced prompt designs should be standard; and
the empty-input probe (does the model compile from just ``Proposition N''?) is a
cheap memorization check worth running before reporting autoformalization
numbers.

\section*{Limits}

\begin{itemize}
\setlength{\itemsep}{0.2em}
\item \textbf{One model.} The counterbalanced results are GPT-4o only; the Qwen
runs are void due to contamination. Prompt sensitivity may differ elsewhere.
\item \textbf{Compiling is not correctness.} Compile rates overstate real
performance. On the subset semantically scored: 14\% strict equivalence, 30\%
correct-direction. The 24-point prompt effect is measured on compilation, not
correctness.
\item \textbf{Uneven cells.} Fitzpatrick-prompt cells carry 430 runs each;
Heath and Modern prompt cells have 133--235. The model handles the imbalance,
but precision varies.
\item \textbf{UniGeo, one prompt.} The 100-theorem replication used only
Fitzpatrick-style prompts, so it confirms the translation non-effect but not the
prompt effect.
\end{itemize}

\vspace{0.6em}
\noindent\rule{\linewidth}{0.4pt}\\[0.3em]
{\small Built on \href{https://github.com/Zed-Rez/youklid}{Youklid} and the
\href{https://github.com/loganrjmurphy/LeanEuclid}{LeanEuclid} benchmark
(Murphy et al., 2024). Contamination probe ran locally on a Jetson Orin Nano
Super.}

\end{document}
